\subsection{Premodular formulas}
\label{subsec:premodular-formulas}

For $n\in\mathbb Z$, define the shifted Gauss sums
\[
  \tau_\eta^+(n)
  :=\sum_{i\in\operatorname{Irr}(\cC)}d_i^2\theta_i\eta_i^n,
  \qquad
  \tau_\eta^-(n)
  :=\sum_{i\in\operatorname{Irr}(\cC)}d_i^2\theta_i^{-1}\eta_i^n.
\]

\begin{proposition}[Odd premodular formula]
\label{prop:odd-formula}
Let $X$ be a closed, oriented, simply connected smooth four-manifold with
odd intersection form.  Choose an integral diagonalization
\[
  Q_X\cong p\langle1\rangle\oplus q\langle-1\rangle
\]
with basis $e_1^+,\ldots,e_p^+,e_1^-,\ldots,e_q^-$, and write
\[
  \alpha=\sum_{r=1}^p a_re_r^++\sum_{s=1}^q b_se_s^-.
\]
Then
\[
  \widehat{\CY}_{\cC,\eta}(X,\alpha)
  =\prod_{r=1}^{p}\tau_\eta^+(a_r)
   \prod_{s=1}^{q}\tau_\eta^-(-b_s).
\]
\end{proposition}

\begin{proof}
The diagonal lattice is represented by the connected sum of $p$ positively
framed and $q$ negatively framed unknots.  The evaluator exponent on a
$2$-handle is the value of $\operatorname{PD}(\alpha)$ on its oriented
generator.  Hence
\[
  \langle\operatorname{PD}(\alpha),e_r^+\rangle=a_r,
  \qquad
  \langle\operatorname{PD}(\alpha),e_s^-\rangle=-b_s.
\]
The corresponding one-component evaluations are therefore
$\tau_\eta^+(a_r)$ and $\tau_\eta^-(-b_s)$.  Apply normalized connected-sum
multiplicativity.  If $X$ is not itself the standard connected sum, use
\cref{cor:marked-stable-realization} and cancel the common nonzero
hyperbolic stabilization factor by
\cref{prop:marked-stabilization-formula}.
\end{proof}

Fix a hyperbolic basis $e,f$ of $H_2(H;\mathbb Z)$, with
$Q_H(e,f)=1$, and write $\alpha=ae+bf$.  Put
\[
  \mathcal H_{\cC,\eta}(a,b)
  :=\sum_{i,j\in\operatorname{Irr}(\cC)}
       d_i d_j\eta_i^b\eta_j^a
       \operatorname{Hopf}_{\cC}^{+}(i,j),
\]
where $\operatorname{Hopf}_{\cC}^{+}(i,j)$ is the positive Hopf-link
scalar, with its first and second components representing $e$ and $f$,
respectively.  The exchanged exponents occur because
$\langle\operatorname{PD}(ae+bf),e\rangle=b$ and
$\langle\operatorname{PD}(ae+bf),f\rangle=a$.
For a marked standard block $(B,\beta)$ define
\[
  V_{\cC,\eta}(B,\beta)=
  \begin{cases}
    \widehat{\CY}_{\cC,\eta}(K3,\beta),&B=K3,\\
    \widehat{\CY}_{\cC,\eta}(\overline{K3},\beta),
      &B=\overline{K3},\\
    \mathcal H_{\cC,\eta}(a,b),
      &B=H\text{ and }\beta=(a,b).
  \end{cases}
\]

\begin{proposition}[Even premodular formula]
\label{prop:even-formula}
Let $X$ be a closed, oriented, simply connected smooth four-manifold with
even intersection form, and let $\alpha\in H_2(X;\mathbb Z)$.  Suppose that
for some $s\geq0$ there is an isometry of marked lattices
\[
  \bigl(H_2(X\#sH;\mathbb Z),Q_{X\#sH},\alpha\oplus0\bigr)
  \cong
  \bigoplus_{r=1}^{N}
    \bigl(H_2(B_r;\mathbb Z),Q_{B_r},\beta_r\bigr),
\]
where every $B_r$ is $K3$, $\overline{K3}$, or $H$.  Then
\[
  \widehat{\CY}_{\cC,\eta}(X,\alpha)
  =\frac{\displaystyle\prod_{r=1}^{N}
      V_{\cC,\eta}(B_r,\beta_r)}
    {\mathcal H_{\cC,\eta}(0,0)^s}.
\]
\end{proposition}

\begin{proof}
By \cref{thm:freedman-donaldson}, the stated marked lattice isometry is
realized by an orientation-preserving homeomorphism from $X\#sH$ to the
connected sum of the $B_r$, carrying $\alpha\oplus0$ to the connected sum
of the $\beta_r$.  This marked homeomorphism is realized after a common
further hyperbolic stabilization by
\cref{cor:marked-stable-realization}.  Apply
\cref{prop:marked-stabilization-formula} to both sides.  Normalized
connected-sum multiplicativity splits the standard side into the displayed
$K3$, $\overline{K3}$, and hyperbolic factors.  The common additional
stabilizations cancel.  Since
$\mathcal H_{\cC,\eta}(0,0)=h_{\cC,\eta}\neq0$, the original $s$
hyperbolic factors may be divided out, giving the formula.
\end{proof}

\begin{corollary}[Simply connected no-go theorem]
\label{cor:no-simply-connected-exotic-detection}
Let $X$ and $Y$ be homeomorphic closed, oriented, simply connected smooth
four-manifolds, and let $f\colon X\to Y$ be an orientation-preserving
homeomorphism.  Then, for every $\alpha\in H_2(X;\mathbb Z)$,
\[
  \CY_{\cC,\eta}(X,\alpha)
  =\CY_{\cC,\eta}(Y,f_*\alpha).
\]
Thus the marked invariant does not distinguish exotic smooth structures on
a simply connected topological four-manifold.
\end{corollary}

\begin{proof}
By \cref{cor:marked-stable-realization}, for some $r$ there is a
diffeomorphism
\[
  (X\#rH,\alpha\oplus0)
  \cong(Y\#rH,f_*\alpha\oplus0).
\]
Marked diffeomorphism invariance identifies the stabilized values.  Applying
\cref{prop:marked-stabilization-formula} to both sides gives a common factor
$h_{\cC,\eta}^{r}$.  Since this factor is nonzero, it cancels and yields
the asserted equality.
\end{proof}
